% Part: conditional-logics
% Chapter: minimal-change-semantics
% Section: true-false

\documentclass[../../../include/open-logic-section]{subfiles}

\begin{document}

\olfileid{con}{min}{tf}

\olsection{प्रतिवास्तविक विधानांचे सत्य आणि असत्य}

सर्वाधिक जवळची सर्व $!A$-जगे
$!B$-जगे असतील, तर प्रतिवास्तविक
विधान $!A \cif !B$
(रिक्त नसलेल्या अर्थाने) सत्य असते.
\olref{fig:true} मध्ये हे दाखवले आहे.
\begin{figure}
\begin{center}
\begin{tikzpicture}[scale=.7]\small
  \spheresystem{5}
  \draw (0,0) node {$w$};
  \propositionintersect{0}{4}{40}{3.5}
  {\tikzset{proposition/.append={smooth,tension=1.4}}
  \proposition[shift={(1.6,-1)}]{90}{1}{30}{4}}
  \path (1.5,3) node[below] {$\formula{B}$};
  \path (3,0) node {$\formula{A}$};
\end{tikzpicture}
\caption{रिक्त नसलेल्या अर्थाने सत्य प्रतिवास्तविक विधान}
\ollabel{fig:true}
\end{center}
\end{figure}
जग~$w$ भोवतीच्या गोलकांच्या
प्रणालीत $!A$ समाविष्ट
करणारा एकही गोलक नसेल,
तरी~$w$ येथे प्रतिवास्तविक
विधान सत्य असते. अशा वेळी
ते \emph{रिक्तपणे} सत्य
असते (\olref{fig:vacuous} पाहा).
\begin{figure}
\begin{center}
\begin{tikzpicture}[scale=.7]\small
  \spheresystem{5}
  \draw (0,0) node {$w$};
  \proposition{0}{6.5}{40}{3.5}
  {\tikzset{proposition/.append={smooth,tension=1.4}}
  \proposition[shift={(1.2,-1)}]{90}{1}{30}{4}}
  \path (1.5,3) node[below] {$\formula{B}$};
  \spherepos{0}{7}{node {$\formula{A}$}}
\end{tikzpicture}
\caption{रिक्तपणे सत्य प्रतिवास्तविक विधान}
\ollabel{fig:vacuous}
\end{center}
\end{figure}

प्रतिवास्तविक विधान दोन
प्रकारे असत्य असू शकते.
पहिल्या प्रकारात सर्वाधिक
जवळची $!A$-जगे
सर्वच $!B$-जगे नसतात,
पण त्यांपैकी काही
तशी असतात.
या वेळी $!A \cif
\lnot !B$ हे विधानही
असत्य असते
(\olref{fig:false} पाहा).
\begin{figure}
\begin{center}
\begin{tikzpicture}[scale=.7]\small
  \spheresystem{5}
  \draw (0,0) node {$w$};
  \propositionintersect{0}{4}{40}{3.5}
  {\tikzset{proposition/.append={smooth,tension=1.4}}
  \proposition[shift={(1.6,0)}]{90}{1}{30}{3}}
  \path (1.5,3) node[below] {$\formula{B}$};
  \path (3,0) node {$\formula{A}$};
\end{tikzpicture}
\caption{असत्य प्रतिवास्तविक विधान; विरुद्ध उत्तरांगाचे विधानही असत्य}
\ollabel{fig:false}
\end{center}
\end{figure}
सर्वाधिक जवळच्या
$!A$-जगांचा
$!B$-जगांशी
अजिबात छेद नसेल,
तर $!A \cif !B$
असत्य असते. पण
या वेळी सर्वाधिक
जवळची सर्व
$!A$-जगे
$\lnot !B$-जगे
असतात; म्हणून
$!A \cif \lnot !B$
सत्य असते
(\olref{fig:false-opposite} पाहा).
\begin{figure}
\begin{center}
\begin{tikzpicture}[scale=.7]\small
  \spheresystem{5}
  \draw (0,0) node {$w$};
  \propositionintersect{0}{4}{40}{3.5}
  {\tikzset{proposition/.append={smooth,tension=1.4}}
  \proposition[shift={(-1,0)}]{90}{1}{30}{3}}
  \path (-1,3) node[below] {$\formula{B}$};
  \path (3,0) node {$\formula{A}$};
\end{tikzpicture}
\caption{असत्य प्रतिवास्तविक विधान; विरुद्ध उत्तरांगाचे विधान सत्य}
\ollabel{fig:false-opposite}
\end{center}
\end{figure}

कठोर अभिव्यंजनाच्या
उलट, प्रतिवास्तविक
विधाने परिस्थितीवश
असू शकतात.
\olref{fig:contingent} मधील
गोलक प्रतिमान पाहा.
$u$ च्या सर्वाधिक
जवळची $!A$-जगे
ही सर्व $!B$-जगे
आहेत; म्हणून
$\mSat{M}{!A \cif
  !B}[u]$. पण~$v$
च्या सर्वाधिक जवळच्या
काही $!A$-जगांत
$!B$ सत्य नाही;
म्हणून
$\mSat/{M}{!A \cif !B}[v]$.
\begin{figure}
\begin{center}
\begin{tikzpicture}[scale=.6]\small
  \spheresystem{7}
  \spheresystem[shift={(0:3.1)},dashed]{4}
  \propositionintersect{45}{4}{40}{4.5}
  \begin{scope}
    \clip \propositionplot{45}{4}{40}{4.5} ;
    \spherelayer[shift={(0:3.1)}]{3}
  \end{scope}
  \proposition[shift={(1.4,-.5)}]{90}{1}{35}{5}
  \path (0,0) node {$u$};
  \path (0:3.1) node {$v$};
  \spherepos{35}{9}{node {$\formula{A}$}}
  \spherepos{80}{9}{node {$\formula{B}$}}
\end{tikzpicture}
\end{center}
\caption{परिस्थितीवश प्रतिवास्तविक विधान}
\ollabel{fig:contingent}
\end{figure}
\end{document}
